% SPDX-License-Identifier: CC-BY-4.0
\documentclass[11pt]{article}

\usepackage[T1]{fontenc}
\usepackage[utf8]{inputenc}
\usepackage{lmodern}
\usepackage{amsmath,amssymb,amsthm,mathtools}
\usepackage{booktabs}
\usepackage[margin=1in]{geometry}
\usepackage{microtype}
\usepackage[colorlinks=true,linkcolor=blue,citecolor=blue,urlcolor=blue]{hyperref}

\hypersetup{
  pdftitle={Cores of the Merged Johnson Graphs J(2k+1,k) with relations 2 through k-2},
  pdfauthor={Jihun Kim},
  pdfkeywords={graph core, graph endomorphism, merged Johnson graph, Johnson scheme, qualitative independence}
}

\theoremstyle{plain}
\newtheorem{theorem}{Theorem}[section]
\newtheorem{lemma}[theorem]{Lemma}
\newtheorem{proposition}[theorem]{Proposition}
\newtheorem{corollary}[theorem]{Corollary}
\theoremstyle{remark}
\newtheorem{remark}[theorem]{Remark}

\newcommand{\Aut}{\operatorname{Aut}}
\newcommand{\End}{\operatorname{End}}
\newcommand{\one}{\boldsymbol{1}}
\newcommand{\supp}{\operatorname{supp}}
\newcommand{\R}{\mathbb R}

\title{Cores of the Merged Johnson Graphs\\
$J(2k+1,k)_{\{2,\ldots,k-2\}}$}
\author{Jihun Kim\thanks{Email:
\href{mailto:jihunkimkw@gmail.com}{\texttt{jihunkimkw@gmail.com}}}}
\date{25 August 2026}

\begin{document}

\maketitle

\begin{abstract}
For $k\ge4$, let $G_k$ be the graph on the $k$-subsets of $[2k+1]$, where
two distinct vertices are adjacent when their intersection has size between
$2$ and $k-2$.  We prove that
\[
 \End(G_k)=\Aut(G_k)\cong S_{2k+1},
\]
and hence that $G_k$ is a core.  We also show that
$\alpha(G_k)=2k+3$ and classify the maximum independent sets.  The proof
reduces a nontrivial core retraction to a translate-transversal identity and
then uses the multiplicity-free Johnson decomposition to separate the
spectral supports of a fibre and the core image.  The resulting support
restrictions would make one of the remaining Johnson relation counts
negative, a contradiction.  It follows that the almost-uniform
$3$-qualitative independence hypergraph
$3\text{-}\mathrm{AUQI}(n,2)$ is a core for every odd $n\ge9$.
\end{abstract}

\noindent\textbf{Keywords.}
Graph core; graph endomorphism; merged Johnson graph; Johnson association
scheme; qualitative independence; extremal set system.

\medskip
\noindent\textbf{MSC 2020.}
05C60, 05C25, 05E30, 05D05.

\section{Introduction}

A finite graph is a \emph{core} if every endomorphism is an automorphism;
equivalently, it has no retraction onto a proper subgraph homomorphically
equivalent to it \cite{HellNesetril1992,HahnTardif1997}.  Determining cores
is particularly subtle for highly symmetric graphs: the automorphism group
describes the invertible maps but does not, by itself, exclude
noninvertible foldings.  General results for symmetric and geometric graphs
give strong restrictions on their cores
\cite{CameronKazanidis2008,GodsilRoyle2011,Roberson2013}, but do not decide
the merged relations considered here.

For $k\ge4$, define
\begin{equation}\label{eq:def-Gk}
 \Omega=\binom{[2k+1]}k,\qquad
 V(G_k)=\Omega,\qquad
 A\sim B\iff 2\le |A\cap B|\le k-2.
\end{equation}
If $d(A,B)=k-|A\cap B|$ is Johnson distance, then the edge set is the
union of relations $2,\ldots,k-2$.  The same interval occurs in both
intersection and distance notation; thus
\[
 G_k=J(2k+1,k)_{\{2,\ldots,k-2\}}.
\]
Jones's classification of
automorphisms of merged Johnson graphs gives
\begin{equation}\label{eq:aut}
 \Aut(G_k)\cong S_{2k+1}
\end{equation}
in its natural action on $\Omega$ \cite[Theorem~2]{Jones2005}.  Indeed, for
the Johnson distance set $I=\{2,\ldots,k-2\}$, the relevant exceptional condition
$I=k+1-I$ does not occur, since
$k+1-I=\{3,\ldots,k-1\}\ne I$.  The remaining question for this family is
whether $G_k$ admits a noninjective endomorphism.

The family arises in qualitative independence.  For odd $n=2k+1$, the
vertices of $3\text{-}\mathrm{AUQI}(n,2)$ are almost-balanced binary
strings, represented by $k$-subsets of $[n]$, and a hyperedge is a triple
realizing all eight binary patterns.  Thomas and Akhtar identify $G_k$ as
the relevant $2$-section for
$3\text{-}\mathrm{AUQI}(2k+1,2)$ and show that if this merged Johnson graph
is a core, then so is the hypergraph
\cite[Lemma~7 and Proposition~7]{ThomasAkhtar2026}.  They pose the odd
cases as a direction for further work.  Theorem~\ref{thm:main} resolves
that question for every odd $n\ge9$.

\begin{theorem}\label{thm:main}
For every integer $k\ge4$,
\[
 \End(G_k)=\Aut(G_k)\cong S_{2k+1}.
\]
In particular, $G_k$ is a core.
\end{theorem}

The proof begins with the extremal and divisibility constraints on the
fibres of a minimum-rank retraction.  Each fibre meets every translate of
the core image in exactly one vertex.  In the language of permutation
groups, this is a non-separation witness for the natural
$S_{2k+1}$-action; see \cite{AljohaniBambergCameron2017}.  The additional
fact that the fibre is independent restricts its distance distribution to
three Johnson relations.  The transversal identity separates the
nonconstant Johnson supports of the fibre and the core image, after which
positivity excludes first the proper divisors of $2k+3$ and then the
remaining fibre size $2k+3$.

We also characterize the maximum independent sets, although only the value
$\alpha(G_k)=2k+3$ is used in the core argument.

\section{The extremal independent families}\label{sec:alpha}

An independent family $\mathcal F\subseteq \Omega$ has
\begin{equation}\label{eq:allowed-intersections}
 |A\cap B|\in\{0,1,k-1\}\qquad(A\ne B\in\mathcal F).
\end{equation}
General restricted-intersection results are given in
\cite{DezaErdosFrankl1978}.  In the case $n=2k+1$ and
$L=\{0,1,k-1\}$, the following proposition gives the sharp bound and
characterizes equality.

\begin{proposition}\label{prop:alpha}
For every $k\ge4$,
\[
 \alpha(G_k)=2k+3.
\]
Moreover, every maximum independent family is, up to a permutation of
$[2k+1]$, of the following form.  Choose disjoint sets $U,T,\{p\}$ with
$|U|=k+1$ and $|T|=k-1$; then
\begin{equation}\label{eq:max-independent}
 \mathcal M(U,T,p)=
 \binom Uk\ \cup\ \{T\cup\{x\}:x\in U\cup\{p\}\}.
\end{equation}
Consequently the number of maximum independent families is
\[
 (k+2)\binom{2k+1}{k-1}.
\]
\end{proposition}

\begin{proof}
Define $A\equiv B$ when $A=B$ or $|A\cap B|=k-1$.  If
$A\equiv B\equiv C$, then $|A\cap C|\ge k-2\ge2$; hence
\eqref{eq:allowed-intersections} shows that $A\equiv C$.  Thus the
$\equiv$-classes are cliques in the ordinary Johnson graph.  Each such
clique is contained in one of the two familiar types
\begin{align*}
 \text{star: }&\quad \{T\cup\{x\}:x\in R\},\qquad
 |T|=k-1,\ R\subseteq[2k+1]\setminus T,\\
 \text{top: }&\quad \{U\setminus\{u\}:u\in R\},\qquad
 |U|=k+1,\ R\subseteq U.
\end{align*}
Indeed, after fixing two members with a common $(k-1)$-set, every third
member either contains that common set or is contained in their
$(k+1)$-element union.

Suppose first that a star component has size at least four.  An outside
block must be disjoint from its common $(k-1)$-set $T$.  Otherwise it
contains at most one point of $T$ and must avoid every point indexing the
star component.  If that component has size $r\ge4$, at most
$1+(k+2-r)<k$ points remain available.
Therefore all outside blocks lie in the $(k+2)$-set
$W=[2k+1]\setminus T$.  Two distinct outside blocks meet in at least
$k-2$ points, while independence restricts their intersection to
$\{0,1,k-1\}$.  Since $k\ge4$, they must meet in $k-1$ points.  Thus all
outside blocks form one Johnson clique of size at most $k+1$.  The original
star has size at most $k+2$, so
$|\mathcal F|\le2k+3$.

Suppose next that a top component has size at least four, contained in
$\binom Uk$ with $|U|=k+1$.  Put $W=[2k+1]\setminus U$, so $|W|=k$.
For an outside block $B$, write $a=|B\cap U|$.  Comparison with
$U\setminus\{u\}$ gives $a-\one_{\{u\in B\}}\le1$ for each index $u$ of
the component.  If $a\ge3$, then
$a-\one_{\{u\in B\}}\ge2$, a contradiction.  If $a=2$, the inequality
forces all four component indices to lie in $B$, contradicting
$|B\cap U|=2$.  Hence $a\le1$, and $B$ is either $W$ or has the form
\[
 \{u\}\cup(W\setminus\{w\}),\qquad u\in U,\quad w\in W.
\]
Two blocks of the latter type are compatible only when they have the same
$u$ or the same $w$.  Thus they lie in a row or a column of the complete
bipartite grid $U\times W$; with $W$ included, their number is at most
$k+2$.  The top has size at most $k+1$, again giving $2k+3$.

It remains to treat families whose components all have size at most three.
Distinct components have disjoint $2$-shadows, since a common pair would
give intersection at least two.  A component of size $r=1,2,3$ contains at
least, respectively,
\[
 \binom{k}{2},\qquad \frac{(k-1)(k+2)}2,\qquad
 \frac{k(k+1)}2
\]
distinct pairs.  Each quantity is at least
$r k(k+1)/6$.  Summing disjoint shadows gives
\[
 |\mathcal F|\frac{k(k+1)}6
 \le \binom{2k+1}{2}=k(2k+1),
\]
and therefore
\[
 |\mathcal F|\le \frac{6(2k+1)}{k+1}<2k+3.
\]
The strict inequality follows from
\[
 (2k+3)(k+1)-6(2k+1)=2k^2-7k-3>0
 \qquad(k\ge4).
\]

The family in \eqref{eq:max-independent} has size $(k+1)+(k+2)$ and
satisfies \eqref{eq:allowed-intersections}, proving the value of
$\alpha(G_k)$.  Equality in either large-component case forces both
component bounds to be tight.  In the star case, the outside $k$-sets in
the $(k+2)$-set $W$ correspond, by complementation, to a maximum
pairwise-intersecting family of $2$-subsets of $W$.  Since $|W|=k+2\ge6$,
a maximum such family is a full star, and
hence a full top on a $(k+1)$-subset $U\subset W$.  The one remaining
point is $p$.  In the top case, equality forces the full top and $k+2$
outside blocks.  A row of $U\times W$, together with $W$, has only $k+1$
blocks, whereas a full column, together with $W$, has $k+2$; the latter is
the full star with common $(k-1)$-set $T=W\setminus\{p\}$.  Thus
\eqref{eq:max-independent} is the only equality type.  Its two
$\equiv$-components have the distinct sizes $k+2$ and $k+1$, so the common
set $T$ and the remaining point $p$ are uniquely recoverable from the
family.  This gives $(k+2)\binom{2k+1}{k-1}$ families.
\end{proof}

\section{From an endomorphism to equal fibres}\label{sec:reduction}

Put
\[
 N=|\Omega|=\binom{2k+1}{k}.
\]
The following lemmas reduce a noninjective endomorphism to two divisibility
conditions and a transversal identity.

\begin{lemma}[Minimum-rank retraction]\label{lem:retraction}
If $G_k$ has a non-automorphic endomorphism, then it has a retraction
$r:G_k\to H$ onto a proper core $H$.  All fibres of $r$ have a common size
$s>1$.  If $q=|H|$, then
\begin{equation}\label{eq:sqN}
 sq=N,\qquad s\mid N.
\end{equation}
Moreover, for every fibre $F_y=r^{-1}(y)$ and every
$\gamma\in S_{2k+1}$,
\begin{equation}\label{eq:transversal}
 |F_y\cap\gamma H|=1.
\end{equation}
\end{lemma}

\begin{proof}
Choose an endomorphism $\phi$ of minimum rank and put
$H=G_k[\phi(\Omega)]$, the subgraph induced by its vertex image.  The restriction
$\phi|_H$ is injective, for otherwise
$\phi^2$ has smaller rank.  Since $H$ is finite, this restriction is an
automorphism.  Composing $\phi$ with its inverse on $H$ gives a retraction
\[
 r=(\phi|_H)^{-1}\circ\phi:G_k\to H.
\]
If $H$ had a noninjective endomorphism, composing it with $r$ would produce
an endomorphism of $G_k$ with smaller rank.  Thus $H$ is a core.

For $y\in H$ and $\gamma\in S_{2k+1}$, the map
$r\circ\gamma|_H:H\to H$ is an endomorphism of the core $H$, hence an
automorphism.  Exactly one point of $\gamma H$ therefore lies in the fibre
$F_y$, proving \eqref{eq:transversal}.

Average \eqref{eq:transversal} over $\gamma$.  Since the point-permutation
group is transitive on $\Omega$, each vertex belongs to exactly
$|S_{2k+1}|q/N$ translates $\gamma H$.  Hence
\[
 |S_{2k+1}|=
 \sum_\gamma |F_y\cap\gamma H|
 =|F_y|\frac{|S_{2k+1}|q}{N}.
\]
Thus $|F_y|=N/q$, independently of $y$, and \eqref{eq:sqN} follows.  This
also gives a direct proof, in the present setting, of the equal-fibre
phenomenon for vertex-transitive graphs
\cite{HahnTardif1997,Roberson2013}.
\end{proof}

\begin{lemma}[Divisibility by \texorpdfstring{$2k+3$}{2k+3}]
\label{lem:second-divisor}
Under the hypotheses of Lemma~\ref{lem:retraction},
\begin{equation}\label{eq:alpha-product}
 s\alpha(H)=2k+3.
\end{equation}
In particular,
\begin{equation}\label{eq:s-divides-m}
 s\mid 2k+3.
\end{equation}
\end{lemma}

\begin{proof}
If $I$ is independent in $H$, then $r^{-1}(I)$ is independent in $G_k$;
equal fibres and Proposition~\ref{prop:alpha} give
\begin{equation}\label{eq:alpha-upper}
 \alpha(H)\le\frac{2k+3}{s}.
\end{equation}
The inclusion $H\hookrightarrow G_k$ and the retraction give homomorphisms
in both directions, so $\chi_f(H)=\chi_f(G_k)$.  A vertex-transitive graph
$G$ satisfies $\chi_f(G)=|V(G)|/\alpha(G)$, while every graph $Y$ satisfies
$\chi_f(Y)\ge |V(Y)|/\alpha(Y)$
\cite[Chapters~1--3]{ScheinermanUllman1997}.  Consequently
\[
 \frac{N}{2k+3}=\chi_f(G_k)=\chi_f(H)
 \ge\frac{q}{\alpha(H)}=\frac{N/s}{\alpha(H)}.
\]
This reverses \eqref{eq:alpha-upper}, proving
\eqref{eq:alpha-product} and hence \eqref{eq:s-divides-m}.
\end{proof}

\section{Spectral transversals in the Johnson scheme}\label{sec:spectral}

Let $A_i$ be the matrix of Johnson distance $i$, so
$A_i(A,B)=1$ precisely when $d(A,B)=i$.  Let $\rho$ denote the permutation
representation of $S_{2k+1}$ on $\R^\Omega$.  Its multiplicity-free
decomposition is
\[
 \R^\Omega=V_0\oplus V_1\oplus\cdots\oplus V_k,
\]
where $V_j\cong S^{(2k+1-j,j)}$.  Write
$\rho_j=\rho|_{V_j}$ and let $E_j$ be orthogonal projection onto $V_j$.
This is the standard Johnson association scheme;
see, for example, \cite{BrouwerCohenNeumaier1989,GodsilNewman2006}.

\begin{lemma}[Spectral-transversal lemma]\label{lem:ST}
Let $F$ be a fibre and let $H$ be the core image in
Lemma~\ref{lem:retraction}.  For every $j\ge1$,
\begin{equation}\label{eq:SC}
 \|E_j\one_F\|^2\,\|E_j\one_H\|^2=0.
\end{equation}
\end{lemma}

\begin{proof}
By \eqref{eq:transversal}, for all $\gamma\in S_{2k+1}$,
\[
 \langle\one_F,\rho(\gamma)\one_H\rangle=1.
\]
As $|F||H|=sq=N$, removing the constant components gives
\[
 \left\langle
 \one_F-\frac{s}{N}\one,
 \rho(\gamma)\left(\one_H-\frac{q}{N}\one\right)
 \right\rangle=0.
\]
After decomposing into $V_1,\ldots,V_k$, we may complexify; the real Specht
modules here are absolutely irreducible.  Schur orthogonality then says that
the matrix-coefficient spaces belonging to distinct irreducibles are
orthogonal.  Since their sum is the zero function, every individual function
\[
 \gamma\longmapsto
 \langle E_j\one_F,\rho_j(\gamma)E_j\one_H\rangle
\]
vanishes identically.  If $E_j\one_H\ne0$, its orbit spans the irreducible
module $V_j$, forcing $E_j\one_F=0$.  This is exactly
\eqref{eq:SC}.
\end{proof}

\subsection{Endpoint eigenvalues and fibre moments}

Use the first-eigenmatrix convention
\[
 A_iE_j=P_j(i)E_j.
\]
The source \cite[Section~5.1 and Theorem~5.4]
{AljohaniBambergCameron2017} uses the opposite order for the two indices.
With our convention, the Eberlein polynomials are
\begin{equation}\label{eq:Eberlein}
 P_j(i)=\sum_{r=0}^{i}(-1)^{i-r}
 \binom{k-r}{i-r}\binom{k-j}{r}\binom{k+1+r-j}{r}.
\end{equation}
Put
\begin{equation}\label{eq:Kteps}
 K=k+1,\qquad t=K-j,\qquad \epsilon=(-1)^j.
\end{equation}
Specializing \eqref{eq:Eberlein}, and using
$A_1A_k=2A_{k-1}+(K-1)A_k$ for the middle endpoint, gives
\begin{equation}\label{eq:endpoint-eigenvalues}
 \boxed{
 P_j(1)=t^2-K,\quad
 P_j(k-1)=\epsilon\frac{t(t^2-2K+1)}2,\quad
 P_j(k)=\epsilon t.}
\end{equation}
Let $v_i$ denote the valency of relation $i$.  At the three relevant
distances,
\begin{equation}\label{eq:valencies}
 v_1=K(K-1),\qquad
 v_{k-1}=\frac{K(K-1)^2}{2},\qquad
 v_k=K.
\end{equation}

Each fibre is independent, so its distinct pairs have only distances
$1,k-1,k$.  Let $x,y,z$ be the numbers of unordered pairs at these three
distances and normalize
\begin{equation}\label{eq:distance-distribution}
 p_1=\frac{2x}{s(s-1)},\qquad
 p_{k-1}=\frac{2y}{s(s-1)},\qquad
 p_k=\frac{2z}{s(s-1)};\qquad p_1+p_{k-1}+p_k=1.
\end{equation}
Here $\one_F^{\mathsf T}A_1\one_F=2x$, and similarly at distances
$k-1$ and $k$; this accounts for the factor $2$.
With $\mu_j=\dim V_j$, the primitive-idempotent expansion
\[
 E_j=\frac{\mu_j}{N}\sum_{i=0}^k\frac{P_j(i)}{v_i}A_i
\]
shows that
\begin{equation}\label{eq:eta-def}
 \|E_j\one_F\|^2=\frac{\mu_j}{N}\eta_j,
\end{equation}
where, writing $e_j=\eta_j/s$,
\begin{equation}\label{eq:e-general}
e_j=1+(s-1)(a_jp_1+b_jp_{k-1}+c_jp_k)
\end{equation}
with
\begin{equation}\label{eq:abc}
 a_j=\frac{t^2-K}{K(K-1)},\quad
 b_j=\frac{\epsilon t(t^2-2K+1)}{K(K-1)^2},\quad
 c_j=\frac{\epsilon t}{K}.
\end{equation}
In particular, every $\eta_j$ and $e_j$ is nonnegative.

For the core image define
\begin{equation}\label{eq:hR}
 h_j=\|E_j\one_H\|^2,\qquad
 R_i=\one_H^{\mathsf T}A_i\one_H
     =\sum_{j=0}^kP_j(i)h_j\ge0.
\end{equation}
Since $sq=N$,
\begin{equation}\label{eq:h-masses}
 h_0=\frac{q^2}{N}=\frac qs,\qquad
 \sum_{j\ge1}h_j=\frac{q(s-1)}s.
\end{equation}
Lemma~\ref{lem:ST} and \eqref{eq:eta-def} give the complementarity
condition
\begin{equation}\label{eq:eta-h}
 \eta_jh_j=0\qquad(1\le j\le k).
\end{equation}

\begin{lemma}[The first module]\label{lem:first-module}
Every fibre satisfies $\eta_1>0$.  Consequently
\begin{equation}\label{eq:h1-zero}
 h_1=0.
\end{equation}
\end{lemma}

\begin{proof}
If $E_1\one_F=0$, then all point degrees in $F$ are equal, with common
value $sk/(2k+1)$.  Since $\gcd(k,2k+1)=1$, integrality gives
$2k+1\mid s$.  But $s\mid2k+3$ by Lemma~\ref{lem:second-divisor}, and
$\gcd(2k+1,2k+3)=1$, a contradiction to $s>1$.  Thus $\eta_1>0$, and
\eqref{eq:h1-zero} follows from \eqref{eq:eta-h}.
\end{proof}

\begin{remark}[A design consequence]\label{rem:design}
For $j=2$, the three normalized endpoint eigenvalues are
\[
 \frac{k-3}{k+1},\qquad
 \frac{(k-1)(k-4)}{k(k+1)},\qquad
 \frac{k-1}{k+1}.
\]
They are nonnegative for $k\ge4$, so the defining moment formula gives
$\eta_2\ge s$ and hence $h_2=0$.  The point--block and pair--block incidence
matrices have
column spaces $V_0\oplus V_1$ and
$V_0\oplus V_1\oplus V_2$, respectively.  Hence vanishing of the first two
nonconstant Johnson components is equivalent to constant point and pair
degrees, that is, to the $2$-design conditions.  Double counting incidences
and point-pairs gives
\[
 r_H=\frac{kq}{2k+1},\qquad
 \lambda_H=\frac{(k-1)q}{2(2k+1)}.
\]
This observation is not needed in the proof of Theorem~\ref{thm:main}.
\end{remark}

\section{Proper divisors of \texorpdfstring{$2k+3$}{2k+3}}\label{sec:proper}

Set
\begin{equation}\label{eq:m}
 m=2k+3=2K+1.
\end{equation}
Suppose $1<s<m$.  Because $s\mid m$ and $m$ is odd,
\begin{equation}\label{eq:small-s-range}
 3\le s\le\frac{2K+1}{3}\le K.
\end{equation}
We now exclude these values of $s$.

\begin{lemma}[Positivity for even modules]\label{lem:T1-even}
If $1\le j\le k$ is even and $3\le s\le K$, then
\begin{equation}\label{eq:T1-even-gap}
 e_j\ge\frac{K-s+1}{K}>0.
\end{equation}
\end{lemma}

\begin{proof}
Here $\epsilon=1$.  The three coefficient shifts are
\begin{align}
 a_j+\frac1K&=\frac{t^2-1}{K(K-1)}\ge0,\label{eq:even-a}\\
 b_j+\frac1K&=
 \frac{t^3-(2K-1)t+(K-1)^2}{K(K-1)^2}>0,\label{eq:even-b}\\
 c_j+\frac1K&=\frac{t+1}{K}>0.\label{eq:even-c}
\end{align}
For the strict sign in \eqref{eq:even-b}, write the numerator as
$\Psi_K(t)$.  If $t=1$, then $\Psi_K(1)=(K-1)(K-3)>0$.  If $t\ge2$, then
$K\ge t+1$, $\Psi_K(t)$ is increasing as a function of integer
$K\ge t+1$, and
$\Psi_{t+1}(t)=t(t^2-t-1)>0$.
Using $p_1+p_{k-1}+p_k=1$ in \eqref{eq:e-general}, we obtain the affine
identity
\begin{equation}\label{eq:T1-even-identity}
 e_j-\frac{K-s+1}{K}
 =(s-1)\left[
 \left(a_j+\frac1K\right)p_1+
 \left(b_j+\frac1K\right)p_{k-1}+
 \left(c_j+\frac1K\right)p_k\right]\ge0.
\end{equation}
\end{proof}

\begin{lemma}[Zero criterion for odd modules]\label{lem:T1-odd}
Suppose $3\le s\le K$ and $1\le j\le k$ is odd.  Put $t=K-j$ and define
\[
 D_j=K-(s-1)t.
\]
If $D_j\ge0$, then $e_j>0$.  Consequently
\begin{equation}\label{eq:zero-criterion}
 e_j=0\quad\Longrightarrow\quad
 j\text{ is odd and }D_j\le-1.
\end{equation}
\end{lemma}

\begin{proof}
Since $\epsilon=-1$,
\begin{equation}\label{eq:AB}
 A_t:=a_j-c_j=\frac{(t-1)(t+K)}{K(K-1)}\ge0,\qquad
 C_t:=b_j-c_j=\frac{t(K^2-t^2)}{K(K-1)^2}>0.
\end{equation}
Eliminating $p_k$ gives
\begin{equation}\label{eq:T1-odd-basic}
 e_j-\frac{D_j}{K}=(s-1)(A_tp_1+C_tp_{k-1}).
\end{equation}
If $D_j>0$, then \eqref{eq:T1-odd-basic} gives
$e_j\ge D_j/K>0$.
\begin{samepage}
Suppose $D_j=0$.  Then
$(s-1)t=K$ and $3\le s\le K$, so
\[
 2\le t\le K/2,\qquad j-1=K-t-1>0.
\]
With
\[
 \lambda_t=\frac{(t-1)(t+K)}{(K-2)(2K-1)},
\]
simplification gives
\begin{equation}\label{eq:T1-D0}
 e_j-\lambda_t e_1-\lambda_t\frac{j-1}{t}
 =\frac{(K+t)(K-t-1)(t(K-2)+1)}
 {t(K-2)(K-1)^2}\,p_{k-1}\ge0.
\end{equation}
The coefficient $\lambda_t$ is chosen to eliminate the $p_1$ term.
Moreover, $\lambda_t(j-1)/t$ is positive.  Since $e_1\ge0$, this
proves $e_j>0$.
\end{samepage}
Together with Lemma~\ref{lem:T1-even}, this yields
\eqref{eq:zero-criterion}; $D_j$ is integral, so $D_j<0$ means
$D_j\le-1$.
\end{proof}

\begin{proposition}\label{prop:proper-impossible}
No proper divisor $s$ of $m$ can be the common fibre size of a proper core
retraction of $G_k$.
\end{proposition}

\begin{proof}
By \eqref{eq:eta-h}, the nonconstant support of $(h_j)$ is contained in
the zero set of $(e_j)$.  Lemmas~\ref{lem:T1-even} and
\ref{lem:T1-odd} show that every supported index is odd and has
$D_j\le-1$.  Let
\[
 W=\sum_{j\ge1}h_j=\frac{q(s-1)}s,\qquad w_j=\frac{h_j}{W}.
\]
For supported odd $j$, \eqref{eq:endpoint-eigenvalues} gives
$P_j(k)=-t$.  Including the constant contribution
$P_0(k)h_0=Kq/s$, we obtain
\begin{equation}\label{eq:Rk-separation}
 \frac{s}{q}R_k
 =K-(s-1)\sum_jw_jt
 =\sum_jw_j\bigl(K-(s-1)t\bigr)
 =\sum_jw_jD_j\le-1.
\end{equation}
Thus $R_k\le-q/s<0$, contradicting the nonnegativity in
\eqref{eq:hR}.
\end{proof}

\section{The case \texorpdfstring{$s=2k+3$}{s=2k+3}}\label{sec:full}

\begin{proposition}\label{prop:full-impossible}
The value $s=2k+3$ cannot be the common fibre size of a proper core
retraction of $G_k$.
\end{proposition}

Assume throughout this section that
\begin{equation}\label{eq:full-s}
 s=m=2K+1.
\end{equation}
For $t=K-j$ and $\epsilon=(-1)^j$, write
$e_{\epsilon,t}:=e_j$.  Equation~\eqref{eq:e-general} becomes the affine
function
\begin{equation}\label{eq:full-bracket}
 e_{\epsilon,t}=1+\frac{2(t^2-K)}{K-1}p_1
 +\frac{2\epsilon t(t^2-2K+1)}{(K-1)^2}p_{k-1}
 +2\epsilon tp_k.
\end{equation}
\begin{samepage}
The required positivity follows from interpolation among three endpoint
functions.
\begin{lemma}[Endpoint interpolation]\label{lem:endpoint-interpolation}
Under the assumption $s=2K+1$, suppose $K\ge5$.
\begin{enumerate}
 \item If $K$ is odd, set
 \[
  B_1=e_{-,K-1}=e_1,\qquad
  B_{k-1}=e_{-,2}=e_{k-1},\qquad
  B_k=e_{+,1}=e_k.
 \]
 On the parity grids
 \[
  (-,t),\ t=2,4,\ldots,K-1,\qquad
  (+,t),\ t=1,3,\ldots,K-2,
 \]
 every $e_{\epsilon,t}$ is a nonnegative linear combination of
 $B_1,B_{k-1},B_k$.  Away from the three displayed endpoints, the
 coefficient of $B_1$ is positive.

 \item If $K$ is even, the same conclusion holds with
 \[
  B_1=e_{-,K-1}=e_1,\qquad
  B_{k-1}=e_{+,2}=e_{k-1},\qquad
  B_k=e_{-,1}=e_k
 \]
 on the grids
 \[
  (-,t),\ t=1,3,\ldots,K-1,\qquad
  (+,t),\ t=2,4,\ldots,K-2.
 \]
\end{enumerate}
\end{lemma}
\end{samepage}

\begin{proof}
Solving the affine $3\times3$ systems gives the factored coefficients in
Appendix~\ref{app:interpolation}.  Their denominators are positive for
$K\ge5$, and the factors in their numerators give the stated signs on the
two parity grids.
\end{proof}

\begin{proof}[Proof of Proposition~\ref{prop:full-impossible}]
First suppose that $K$ is odd.  Since $B_1=e_1>0$ by
Lemma~\ref{lem:first-module}, Lemma~\ref{lem:endpoint-interpolation} shows that
every non-endpoint $e_j$ is positive.  Complementarity
\eqref{eq:eta-h} and $h_1=0$ therefore confine the nonconstant core support
to $j=k-1,k$.

At these endpoints,
\[
\begin{array}{c|cc}
 &P_j(1)&P_j(k)\\ \hline
 j=k-1&4-K&-2\\
 j=k&1-K&1 .
\end{array}
\]
Using $h_{k-1}+h_k=q(s-1)/s$, $s-1=2K$, and the constant contribution
$P_0(1)+P_0(k)=K^2$, we obtain
\begin{equation}\label{eq:odd-core-identity}
 R_1+R_k=\frac qsK(4-K)<0.
\end{equation}
This contradicts $R_1,R_k\ge0$ for every odd $K\ge5$.

Now suppose that $K$ is even.  The same argument again leaves only
$j=k-1,k$.  Here
\[
 P_{k-1}(1)=4-K,\qquad P_k(1)+3=4-K.
\]
Together with $P_0(1)=K(K-1)$ and $s-1=2K$, this gives
\begin{equation}\label{eq:even-core-identity}
 R_1+3h_k=\frac qsK(7-K).
\end{equation}
For even $K\ge8$ the right side is negative and the left side is
nonnegative.  The only smaller even value is $K=6$ ($k=5$), but then
$m=13$ does not divide $N=\binom{11}{5}=462$, so this fibre size is not a
candidate.
\end{proof}

\section{Proof of the main theorem and consequences}

\begin{proof}[Proof of Theorem~\ref{thm:main}]
Suppose that $G_k$ has a non-automorphic endomorphism.  By
Lemma~\ref{lem:retraction}, a minimum-rank choice gives a proper core
retraction with common fibre size $s>1$.  Lemmas~\ref{lem:retraction} and
\ref{lem:second-divisor} give
\[
 s\mid\binom{2k+1}{k},\qquad s\mid 2k+3.
\]
If $s<2k+3$, Proposition~\ref{prop:proper-impossible} gives a
contradiction.  If $s=2k+3$, Proposition~\ref{prop:full-impossible} gives
a contradiction.  Hence every endomorphism is injective.  On a finite
graph it is therefore bijective; a bijective self-homomorphism permutes the
finite edge set injectively and hence is an automorphism.  Finally
\eqref{eq:aut} identifies the group as $S_{2k+1}$.
\end{proof}

\begin{corollary}\label{cor:AUQI}
For every odd integer $n\ge9$, the almost-uniform qualitative independence
hypergraph $3\text{-}\mathrm{AUQI}(n,2)$ is a core.
\end{corollary}

\begin{proof}
Write $n=2k+1$, so $k\ge4$.  Theorem~\ref{thm:main} says that the merged
Johnson graph $G_k$ is a core.  The transfer result of Thomas and Akhtar
\cite[Lemma~7 and Proposition~7]{ThomasAkhtar2026} then gives the stated
hypergraph conclusion.
\end{proof}

\appendix

\section{Coefficients for endpoint interpolation}\label{app:interpolation}

This appendix records the coefficients used in
Lemma~\ref{lem:endpoint-interpolation}.  They are obtained by substituting
$p_k=1-p_1-p_{k-1}$ into \eqref{eq:full-bracket} and solving in the affine
basis $1,p_1,p_{k-1}$; direct substitution verifies the identities below.

\subsection*{Odd \(K\)}

With the endpoint functions from part~(1) of
Lemma~\ref{lem:endpoint-interpolation}, the negative-parity functions satisfy
\begin{equation}\label{eq:odd-minus-interp}
 e_{-,t}=\alpha^-_OB_1+\beta^-_OB_{k-1}+\gamma^-_OB_k,
\end{equation}
where
\begin{align*}
 \alpha^-_O&=\frac{(t-2)(t+1)[K(K-t-1)+4t-3]}
 {K(K-3)(4K-7)},\\
 \beta^-_O&=\frac{(t+1)(K-t-1)[t(2K-5)+3]}
 {3(K-3)(4K-7)},\\
 \gamma^-_O&=\frac{(t-2)(K-t-1)[4t(K-1)+3]}
 {3K(4K-7)}.
\end{align*}
The positive-parity functions satisfy
\begin{equation}\label{eq:odd-plus-interp}
 e_{+,t}=\alpha^+_OB_1+\beta^+_OB_{k-1}+\gamma^+_OB_k,
\end{equation}
where
\begin{align*}
 \alpha^+_O&=\frac{(t-1)(t+2)[K(K-1)+t(K-4)-3]}
 {K(K-3)(4K-7)},\\
 \beta^+_O&=\frac{(t-1)(K+t-1)[t(2K-5)-3]}
 {3(K-3)(4K-7)},\\
 \gamma^+_O&=\frac{(t+2)(K+t-1)[4t(K-1)-3]}
 {3K(4K-7)}.
\end{align*}
For odd $K\ge5$, the denominators are positive.  On the grids in
Lemma~\ref{lem:endpoint-interpolation}, every numerator is nonnegative, and
$\alpha^\pm_O$ is positive away from the three endpoints.

\subsection*{Even \(K\)}

With the endpoint functions from part~(2) of the lemma, the negative-parity
functions satisfy
\begin{equation}\label{eq:even-minus-interp}
 e_{-,t}=\alpha^-_EB_1+\beta^-_EB_{k-1}+\gamma^-_EB_k,
\end{equation}
where
\begin{align*}
 \alpha^-_E&=\frac{(t-1)(t+2)[K(K-t+1)+2t-1]}
 {(K-2)(K+1)(4K-3)},\\
 \beta^-_E&=\frac{(t-1)(K-t-1)[t(2K-1)+1]}
 {3(K+1)(4K-3)},\\
 \gamma^-_E&=\frac{(t+2)(K-t-1)[4t(K-2)+5]}
 {3(K-2)(4K-3)}.
\end{align*}
The positive-parity functions satisfy
\begin{equation}\label{eq:even-plus-interp}
 e_{+,t}=\alpha^+_EB_1+\beta^+_EB_{k-1}+\gamma^+_EB_k,
\end{equation}
where
\begin{align*}
 \alpha^+_E&=\frac{(t-2)(t+1)[K(K+t+1)-2t-1]}
 {(K-2)(K+1)(4K-3)},\\
 \beta^+_E&=\frac{(t+1)(K+t-1)[t(2K-1)-1]}
 {3(K+1)(4K-3)},\\
 \gamma^+_E&=\frac{(t-2)(K+t-1)[4t(K-2)-5]}
 {3(K-2)(4K-3)}.
\end{align*}
For even $K\ge6$, the same sign conclusions hold on the corresponding
parity grids.

\section*{Code availability}

The supplement contains two exact-arithmetic Python scripts.  They check
selected identities symbolically, reconstruct the interpolation
coefficients, and run the finite-range tests listed in the supplementary
README; they do not mechanize the combinatorial or representation-theoretic
arguments.  The scripts were tested with Python~3.9 and SymPy~1.14 and can
be run with \texttt{make check}.  Version~1.0.0 is archived at
\href{https://doi.org/10.5281/zenodo.22096460}{doi:10.5281/zenodo.22096460}.

\section*{Declaration on AI-assisted tools}

The author used OpenAI language models for conjecture exploration, symbolic
calculation, proof review, and prose editing, and reviewed all resulting
material.  The manuscript is licensed under CC BY 4.0; the accompanying
Python scripts are licensed under the MIT License.

\bibliographystyle{plain}
\bibliography{references}

\end{document}
